% Compte rendu du projet Pod Royale (INFO910)
% Compilation : pdflatex rapport.tex (deux passes, pour la table des matières)
\documentclass[11pt,a4paper]{article}

% ---------------------------------------------------------------- Polices, langue
\usepackage[utf8]{inputenc}
\usepackage[T1]{fontenc}
\usepackage[french]{babel}
\usepackage{XCharter}
\usepackage[scaled=0.92]{helvet}
\usepackage[varqu,scaled=0.9]{inconsolata}
\newcommand{\titlefont}{\sffamily\bfseries}

% ---------------------------------------------------------------- Mise en page
\usepackage[a4paper,margin=2.3cm,headheight=14pt]{geometry}
\usepackage[table]{xcolor}
\usepackage{tikz}
\usetikzlibrary{positioning,calc,arrows.meta,fit,backgrounds,shapes.geometric,babel}
\usepackage{fancyvrb}
\usepackage{array,longtable,booktabs,tabularx}
\usepackage{fancyhdr}
\usepackage{microtype}
\usepackage[hidelinks]{hyperref}

\definecolor{kube}{HTML}{326CE5}
\definecolor{nuit}{HTML}{1B2233}
\definecolor{sang}{HTML}{C0392B}
\definecolor{or}{HTML}{D4A017}
\definecolor{gris}{HTML}{6B7280}
\definecolor{fond}{HTML}{F4F6FA}

\hypersetup{
  pdftitle={Pod Royale — Compte rendu INFO910},
  pdfauthor={Palanca \& Grandy},
}

% Titres de sections en sans-serif coloré
\makeatletter
\renewcommand\section{\@startsection{section}{1}{\z@}%
  {-3.2ex \@plus -1ex \@minus -.2ex}{1.6ex \@plus .2ex}%
  {\titlefont\Large\color{nuit}}}
\renewcommand\subsection{\@startsection{subsection}{2}{\z@}%
  {-2.6ex \@plus -1ex \@minus -.2ex}{1.0ex \@plus .2ex}%
  {\titlefont\large\color{kube}}}
\renewcommand\@seccntformat[1]{%
  {\color{sang}\csname the#1\endcsname}\hspace{0.6em}}
\makeatother

% En-têtes et pieds de page
\pagestyle{fancy}
\fancyhf{}
\fancyhead[L]{\small\sffamily\color{gris}Pod Royale}
\fancyhead[R]{\small\sffamily\color{gris}INFO910 — Compte rendu de projet}
\fancyfoot[C]{\small\sffamily\color{gris}\thepage}
\renewcommand{\headrulewidth}{0.4pt}
\renewcommand{\headrule}{{\color{gris!40}\hrule height \headrulewidth}}

\setlength{\parindent}{0pt}
\setlength{\parskip}{0.6em}

% Code
\newcommand{\code}[1]{\texttt{#1}}
\DefineVerbatimEnvironment{Code}{Verbatim}{%
  fontsize=\small, frame=leftline, framerule=2.5pt, rulecolor=\color{kube},
  xleftmargin=6pt, framesep=8pt, formatcom=\color{nuit}}

% Tableaux
\renewcommand{\arraystretch}{1.35}
\newcolumntype{L}[1]{>{\raggedright\arraybackslash}p{#1}}
\newcommand{\entete}[1]{\textcolor{white}{\sffamily\bfseries #1}}

% Encadré
\newcommand{\encadre}[1]{%
  \par\noindent
  \begin{tikzpicture}
    \node[fill=fond, draw=none, inner sep=10pt, text width=\linewidth-24pt,
          align=left] (b) {#1};
    \draw[or, line width=3pt] (b.north west) -- (b.south west);
  \end{tikzpicture}\par}

\begin{document}

% ================================================================ Page de garde
\begin{titlepage}
\begin{tikzpicture}[remember picture, overlay]
  % Bandeau sombre
  \fill[nuit] (current page.north west) rectangle
    ($(current page.north east)+(0,-15.5cm)$);
  \fill[sang] ($(current page.north west)+(0,-15.5cm)$) rectangle
    ($(current page.north east)+(0,-15.8cm)$);

  % L'arène : une grille de pods, les morts barrés, le survivant en or
  \begin{scope}[shift={($(current page.north)+(0,-4.3cm)$)}]
    \foreach \i in {0,...,7} {
      \pgfmathsetmacro{\x}{(\i-3.5)*1.55}
      \node[regular polygon, regular polygon sides=6, minimum size=1.15cm,
            line width=1.2pt, draw=kube!70, fill=kube!15!nuit] (p\i) at (\x,0) {};
    }
    \foreach \i in {0,2,4,6,7} {
      \draw[sang, line width=1.6pt] ($(p\i.center)+(-0.22,-0.22)$) -- ++(0.44,0.44)
                                    ($(p\i.center)+(-0.22,0.22)$) -- ++(0.44,-0.44);
      \node[regular polygon, regular polygon sides=6, minimum size=1.15cm,
            draw=gris, line width=1.2pt] at (p\i) {};
    }
    \foreach \i in {1,5} \fill[kube!80] (p\i.center) circle (0.14);
    \node[regular polygon, regular polygon sides=6, minimum size=1.15cm,
          line width=1.6pt, draw=or, fill=or!25!nuit] at (p3) {};
    \fill[or] ($(p3.center)+(-0.22,-0.12)$) -- ++(0,0.3) -- ++(0.11,-0.15)
      -- ++(0.11,0.22) -- ++(0.11,-0.22) -- ++(0.11,0.15) -- ++(0,-0.3) -- cycle;
  \end{scope}

  % Titre
  \node[anchor=center, text=white] at ($(current page.north)+(0,-7.6cm)$)
    {\titlefont\fontsize{54}{60}\selectfont POD ROYALE};
  \node[anchor=center, text=or] at ($(current page.north)+(0,-9.2cm)$)
    {\sffamily\Large Un battle royale où chaque pod est un guerrier};
  \node[anchor=center, text=white!75, text width=13cm, align=center]
    at ($(current page.north)+(0,-11.4cm)$)
    {\sffamily\normalsize\hyphenpenalty=10000 Les guerriers s'attaquent en appelant le Service
     \texttt{guerrier}~: c'est le load balancing de Kubernetes qui choisit
     la victime au hasard. Le dernier survivant gagne.};

  % Bas de page
  \node[anchor=center, text=nuit] at ($(current page.center)+(0,-4.2cm)$)
    {\titlefont\LARGE Compte rendu de projet};
  \node[anchor=center, text=kube] at ($(current page.center)+(0,-5.3cm)$)
    {\sffamily\Large INFO910 — Kubernetes};
  \draw[gris!50, line width=0.6pt]
    ($(current page.center)+(-3,-6.3cm)$) -- ($(current page.center)+(3,-6.3cm)$);
  \node[anchor=center, text=nuit] at ($(current page.center)+(0,-7.3cm)$)
    {\sffamily\large Palanca \quad\textcolor{sang}{\textbullet}\quad Grandy};
  \node[anchor=center, text=gris] at ($(current page.south)+(0,2.2cm)$)
    {\sffamily Octobre 2026};
\end{tikzpicture}
\end{titlepage}

% ================================================================ Sommaire
\tableofcontents
\vspace{1.5em}

\section*{Le principe}
\addcontentsline{toc}{section}{Le principe}

\textbf{Pod Royale} est un battle royale où chaque pod est un guerrier.
Les guerriers s'attaquent en appelant le Service \code{guerrier}~: c'est le
load balancing de Kubernetes qui choisit la victime au hasard. Les morts
échouent leur \emph{readinessProbe} et sortent des endpoints. Le dernier
survivant gagne, puis un nouveau round commence.

La figure~\ref{fig:apercu} montre l'arène pendant un round~: les cartes des
guerriers avec leurs points de vie, le fil des combats et le tableau d'honneur.

\begin{figure}[h]
\centering
\begin{tikzpicture}[font=\sffamily\footnotesize, x=1cm, y=1cm]
  % Une carte de guerrier : #1 position, #2 nom, #3 pod, #4 PV, #5 PV max, #6 ligne du bas
  \newcommand{\carte}[6]{%
    \begin{scope}[shift={#1}]
      \draw[kube, fill=kube!6, rounded corners=3pt, line width=0.8pt] (0,0) rectangle (4.3,2.1);
      \node[anchor=north west, font=\sffamily\small\bfseries, text=nuit] at (0.15,2.0) {#2};
      \node[anchor=north west, font=\ttfamily\scriptsize, text=gris] at (0.15,1.5) {#3};
      \fill[gris!25] (0.2,0.8) rectangle (3.0,1.0);
      \pgfmathsetmacro{\larg}{2.8*#4/#5}
      \ifdim #4pt>0pt \fill[kube] (0.2,0.8) rectangle ++(\larg,0.2); \fi
      \node[anchor=west, font=\sffamily\scriptsize, text=nuit] at (3.05,0.9) {#4/#5};
      \node[anchor=south west, font=\sffamily\scriptsize, text=gris] at (0.15,0.1) {#6};
    \end{scope}}
  \node[anchor=west, font=\titlefont\normalsize, text=nuit] at (0,5.0) {POD ROYALE};
  \node[anchor=east, text=nuit] at (15.2,5.0)
    {ROUND 7 \quad \textcolor{sang}{\textbf{FIGHTING}} \quad 3\,/\,6 en vie};
  \carte{(0,2.3)}{Gérard le Fourbe}{guerrier-7d9f8b-x2k4p}{98}{130}{2 victimes}
  \carte{(4.6,2.3)}{Josiane la Sans-Pitié}{guerrier-7d9f8b-q8m1z}{41}{100}{1 victime \quad 1 redémarrage}
  \carte{(0,0)}{Kevin l'Inoxydable}{guerrier-7d9f8b-z7c3d}{0}{160}{\textcolor{sang}{\textbf{mort}}}
  \begin{scope}[shift={(4.6,0)}]
    \draw[gris, dashed, rounded corners=3pt, line width=0.8pt] (0,0) rectangle (4.3,2.1);
    \node[font=\sffamily\small\bfseries, text=gris] at (2.15,1.25) {Monique la Chauve};
    \node[text=gris] at (2.15,0.7) {(en attente)};
  \end{scope}
  % Fil des combats et tableau d'honneur
  \begin{scope}[shift={(9.2,0)}]
    \draw[nuit, fill=fond, rounded corners=3pt, line width=0.8pt] (0,1.75) rectangle (6,4.4);
    \node[anchor=north west, font=\titlefont\scriptsize, text=kube] at (0.15,4.3) {FIL DES COMBATS};
    \node[anchor=north west, text width=5.3cm, align=left, text=nuit] at (0.15,3.85) {%
      Josiane a OOMKillé Kevin l'Inoxydable\\
      \textcolor{or!80!black}{\textbf{CRITIQUE~!}} Gérard inflige 28 dégâts\\
      Ginette a été foudroyée par les dieux\\
      ROUND 7 — 6 combattants. FIGHT~!};
    \draw[nuit, fill=fond, rounded corners=3pt, line width=0.8pt] (0,0) rectangle (6,1.55);
    \node[anchor=north west, font=\titlefont\scriptsize, text=or!80!black] at (0.15,1.45) {TABLEAU D'HONNEUR};
    \node[anchor=north west, text width=5.3cm, align=left, text=nuit] at (0.15,1.0) {%
      1. Gérard le Fourbe (Barbare) \hfill \textbf{3}\\
      2. Josiane la Sans-Pitié (Archer) \hfill \textbf{2}};
  \end{scope}
\end{tikzpicture}
\caption{Aperçu de l'arène pendant un round.}
\label{fig:apercu}
\end{figure}

% ================================================================ Architecture
\section{Architecture}

\begin{figure}[h]
\centering
\begin{tikzpicture}[
    font=\sffamily\small,
    >={Stealth[length=2.6mm]},
    boite/.style={draw=#1, fill=#1!8, rounded corners=4pt, line width=1pt,
                  align=center, inner sep=6pt, minimum width=2.9cm},
    pod/.style={regular polygon, regular polygon sides=6, draw=kube,
                fill=kube!12, line width=0.9pt, minimum size=1.05cm, inner sep=0pt},
    lien/.style={->, line width=0.9pt, color=nuit!80},
    note/.style={font=\sffamily\footnotesize, color=gris, align=center},
  ]
  \node[boite=gris] (nav) {\textbf{Navigateur}};
  \node[boite=kube, below=1.1cm of nav] (ing) {\textbf{Ingress}\\
    \footnotesize\texttt{pod-royale.local}};
  \node[boite=kube, below=1.1cm of ing] (arena) {\textbf{arena} (×1)\\
    \footnotesize Flask + UI};
  \node[boite=sang, right=4.2cm of arena] (redis) {\textbf{redis-0}\\
    \footnotesize StatefulSet + PVC};

  \node[pod, below left=1.7cm and 0.4cm of arena] (g1) {g1};
  \node[pod, right=0.35cm of g1] (g2) {g2};
  \node[pod, right=0.35cm of g2] (g3) {g3};
  \node[right=0.15cm of g3, color=kube] (gdots) {\dots};
  \node[draw=kube, dashed, rounded corners=6pt, line width=0.9pt, inner sep=9pt,
        fit=(g1)(gdots), label={[font=\sffamily\footnotesize\bfseries,
        color=kube]above:Deployment guerrier (×N)}] (dep) {};

  \node[boite=kube, below=1.5cm of dep] (svc) {\textbf{Service \texttt{guerrier}}\\
    \footnotesize ClusterIP};

  \draw[lien] (nav) -- node[note, right] {\texttt{http://pod-royale.local}} (ing);
  \draw[lien] (ing) -- (arena);
  \draw[lien, <->] (arena) -- node[note, above] {lit / écrit l'état}
                              node[note, below] {gère les rounds} (redis);
  \draw[lien, <->] (dep.east) -| node[note, pos=0.25, below] {état, PV, events} (redis);

  \foreach \g in {g1,g2,g3}
    \draw[lien, color=sang] (\g.south) -- (\g.south |- svc.north);
  \node[note, color=sang, left=0.2cm of svc.north west, anchor=south east]
    {\texttt{POST /hit}};
  \draw[lien, color=or!90!black, line width=1.1pt]
    (svc.east) -- ++(1.4,0) |- node[note, pos=0.25, right, color=or!70!black]
    {kube-proxy\\choisit la victime\\au hasard} ($(dep.east)+(0,-0.35)$);

  \begin{scope}[on background layer]
    \node[draw=gris!50, fill=fond, rounded corners=8pt, inner sep=14pt,
          fit=(ing)(arena)(redis)(dep)(svc),
          label={[font=\sffamily\footnotesize, color=gris, anchor=south east]
                 south east:namespace \texttt{pod-royale}}] {};
  \end{scope}
\end{tikzpicture}
\caption{Les composants de Pod Royale dans le cluster.}
\end{figure}

\begin{longtable}{L{2.3cm} L{13.2cm}}
\rowcolor{nuit}\entete{Composant} & \entete{Rôle}\\
\endhead
\rowcolor{fond}\code{guerrier/} & Python 3.12 + Flask (waitress). Expose
  \code{/hit}, \code{/ready}, \code{/healthz}, \code{/freeze}, \code{/}.
  Une boucle d'attaque tourne dans un thread.\\
\code{arena/} & Python 3.12 + Flask (waitress). L'arbitre~: il gère les rounds,
  sert la page (Jinja2) et \code{/api/state}. \textbf{1 seul réplica.}\\
\rowcolor{fond}\code{redis} & \code{redis:7-alpine} avec mot de passe et AOF,
  sur un PVC. Stocke l'état du jeu (clés \code{pr:*}). La perte de PV et la
  détection de mort se font dans un script Lua, donc de façon atomique.\\
\bottomrule
\end{longtable}

\encadre{Chaque attaque ouvre une \textbf{nouvelle connexion TCP} (pas de
\code{requests.Session}, en-tête \code{Connection: close}). Sans ça, le
keep-alive collerait chaque attaquant à la même cible et le hasard du load
balancing disparaîtrait.}

% ================================================================ Notions
\clearpage
\section{Notions du cours et où elles apparaissent}

\begin{longtable}{L{3.5cm} L{12.0cm}}
\rowcolor{nuit}\entete{Notion} & \entete{Où dans le projet}\\
\endhead
\rowcolor{fond}\textbf{Pod} & Chaque guerrier est un pod. Son nom de pod donne
  son nom d'affichage («~Gérard le Fourbe~»).\\
\textbf{Namespace} & \code{k8s/00-namespace.yaml}~: tout vit dans
  \code{pod-royale}, et \code{reset.sh} supprime tout d'un coup.\\
\rowcolor{fond}\textbf{Deployment} & \code{k8s/20-guerrier-deployment.yaml}
  (6 réplicas, RollingUpdate \code{maxSurge: 2} / \code{maxUnavailable: 1}) et
  \code{k8s/30-arena-deployment.yaml} (1 réplica, \code{Recreate}).\\
\textbf{Service et DNS} & \code{k8s/21-guerrier-service.yaml}~: la cible des
  attaques, le load balancing choisit la victime.
  \code{redis.pod-royale.svc.cluster.local} est résolu par le DNS du cluster.\\
\rowcolor{fond}\textbf{StatefulSet} & \code{k8s/10-redis-statefulset.yaml}~:
  identité stable \code{redis-0}.\\
\textbf{PVC} & \code{volumeClaimTemplates} de Redis~: le tableau d'honneur
  survit à \code{kubectl delete pod redis-0}.\\
\rowcolor{fond}\textbf{ConfigMap (env)} & \code{jeu-config} via
  \code{envFrom}~: \code{ATTACK\_INTERVAL\_MS}, \code{VICTORY\_PAUSE\_S},
  \code{REDIS\_HOST}…\\
\textbf{ConfigMap (volume)} & La même ConfigMap, montée en
  \code{/config/classes.json}~: les classes de guerriers.\\
\rowcolor{fond}\textbf{Secret} & \code{k8s/01-secret-redis.yaml}~:
  \code{REDIS\_PASSWORD} injecté par \code{secretKeyRef}. Aucun mot de passe en
  clair dans les Deployments.\\
\textbf{readinessProbe} & \code{GET /ready} renvoie 503 quand le guerrier est
  mort, donc il sort des endpoints en $\sim$2\,s (\code{periodSeconds: 2},
  \code{failureThreshold: 1}).\\
\rowcolor{fond}\textbf{livenessProbe} & \code{GET /healthz} renvoie 500 si la
  boucle d'attaque est figée (\code{/freeze}), donc le conteneur redémarre
  (RESTARTS +1).\\
\textbf{startupProbe} & Laisse jusqu'à 60\,s au guerrier pour démarrer avant que
  la livenessProbe prenne le relais~: sur un nœud chargé, un démarrage lent ne
  finit plus en CrashLoopBackOff.\\
\rowcolor{fond}\textbf{requests / limits} & Sur tous les conteneurs (ex.
  guerrier~: \code{20m/48Mi} $\rightarrow$ \code{500m/128Mi}).\\
\textbf{Downward API} & \code{POD\_NAME}, \code{POD\_IP}, \code{NODE\_NAME}
  injectés depuis \code{metadata.name}, \code{status.podIP},
  \code{spec.nodeName}.\\
\rowcolor{fond}\textbf{Arrêt propre (SIGTERM)} &
  \code{terminationGracePeriodSeconds: 10}~: le guerrier annonce «~foudroyé par
  les dieux~» avant de partir.\\
\textbf{Ingress} & \code{k8s/40-ingress.yaml}~: \code{pod-royale.local}
  $\rightarrow$ \code{arena:80}. Seule l'arène est exposée.\\
\rowcolor{fond}\textbf{Gateway API (bonus)} &
  \code{k8s/bonus/41-httproute.yaml}~: \code{Gateway} + \code{HTTPRoute}
  équivalents à l'Ingress. Non appliqués par défaut.\\
\bottomrule
\end{longtable}

% ================================================================ Installation
\section{Installation (minikube)}

Prérequis~: Docker, \code{kubectl}, \code{minikube} (ou \code{kind}).

\subsection{En une commande}

\begin{Code}
./scripts/up.sh                       # macOS, Linux, WSL, Git Bash
\end{Code}

\begin{Code}
powershell -ExecutionPolicy Bypass -File scripts\up.ps1    # Windows
\end{Code}

Le script démarre Docker Desktop s'il le faut (macOS et Windows), puis minikube
si aucun cluster ne répond. Il construit et charge les images, déploie et attend
que tout soit prêt. Il ouvre ensuite l'arène sur
\textbf{\code{http://localhost:8080}} par un port-forward~: pas de \code{sudo},
pas de \code{minikube tunnel}, pas de fichier hosts à modifier, quel que soit
l'OS.

\begin{longtable}{L{4.5cm} L{3.5cm} L{7.0cm}}
\rowcolor{nuit}\entete{\code{up.sh}} & \entete{\code{up.ps1}} & \entete{Effet}\\
\endhead
\rowcolor{fond}\code{PORT=9090} & \code{-Port 9090} & Change le port local.\\
\code{NO\_BROWSER=1} & \code{-NoBrowser} & N'ouvre pas le navigateur.\\
\rowcolor{fond}\code{INGRESS=1} & \code{-Ingress} & Active l'addon ingress pour
  \code{pod-royale.local}.\\
\bottomrule
\end{longtable}

Ctrl+C coupe l'accès, et le jeu continue de tourner dans le cluster.

\subsection{Étape par étape}

\begin{Code}
# 1. Cluster + contrôleur Ingress
minikube start
minikube addons enable ingress

# 2. Construire les images et les charger dans minikube
./scripts/build.sh

# 3. Déployer et attendre que tout soit prêt
./scripts/deploy.sh

# 4. Résoudre pod-royale.local vers le cluster
#    Linux, ou minikube avec un driver VM :
echo "$(minikube ip) pod-royale.local" | sudo tee -a /etc/hosts
#    macOS avec le driver docker : l'IP de minikube n'est pas
#    joignable depuis le Mac. Lancer « minikube tunnel » dans un
#    autre terminal (à laisser ouvert), puis :
echo "127.0.0.1 pod-royale.local" | sudo tee -a /etc/hosts

# 5. Ouvrir http://pod-royale.local
\end{Code}

\textbf{Sans Ingress~:}

\begin{Code}
kubectl -n pod-royale port-forward svc/arena 8080:80
# puis http://localhost:8080
\end{Code}

Pour tout supprimer~: \code{./scripts/reset.sh} (le PVC, et donc le tableau
d'honneur, part avec le namespace).

% ================================================================ Démo
\section{Déroulé de la démo}

\code{./scripts/demo.sh} enchaîne les étapes, avec une pause entre chacune.
L'arène reste ouverte à côté.

\begin{longtable}{>{\centering\arraybackslash}p{0.5cm} L{7.2cm} L{7.2cm}}
\rowcolor{nuit}\entete{\#} & \entete{Commande} & \entete{Ce qu'on voit dans l'arène}\\
\endhead
\rowcolor{fond}1 & \code{kubectl get all -n pod-royale} & Tous les objets du jeu.\\
2 & \code{kubectl get endpoints guerrier -n pod-royale} (pendant un round) &
  Les morts disparaissent des endpoints (readinessProbe).\\
\rowcolor{fond}3 & \code{kubectl scale deploy/guerrier -{}-replicas=12} &
  Des renforts arrivent, en spectateurs jusqu'au prochain round.\\
4 & \code{kubectl delete pod <guerrier>} & «~X a été foudroyé par les
  dieux~», puis un nouveau pod apparaît.\\
\rowcolor{fond}5 & \code{kubectl exec <pod> -{}- curl -X POST
  localhost:8080/freeze} & Le guerrier se fige, la livenessProbe échoue,
  RESTARTS +1 et «~X revient d'entre les morts~».\\
6 & \code{kubectl delete pod redis-0} & Redis redémarre et le tableau
  d'honneur est conservé (PVC).\\
\rowcolor{fond}7 & Patch de la ConfigMap \code{ATTACK\_INTERVAL\_MS=400} +
  \code{kubectl rollout restart deploy/guerrier} & Mode frénésie~: les rounds
  durent quelques secondes.\\
8 & \code{kubectl scale deploy/guerrier -{}-replicas=2} & Le duel final.\\
\bottomrule
\end{longtable}

\textbf{Dans l'arène}, chaque coup réellement échangé entre deux pods est rejoué
par de petits gladiateurs~: flèches pour les archers, boules de feu pour les
mages, entailles à l'épée, à la hache ou à la dague, et le sysadmin qui frappe
au clavier. Un \textbf{coup critique} se voit de loin~: aura dorée pendant la
charge, ralenti, flash et texte «~CRITIQUE~!~». La \textbf{barre d'espace} met
la scène en pause pour expliquer ce qui se passe.

\encadre{Pour voir le hasard du load balancing~:\\[4pt]
\code{kubectl -n pod-royale logs -l app=guerrier -{}-prefix | grep '}$\rightarrow$\code{'}\\[4pt]
Chaque ligne indique quel pod a frappé quel pod.}

% ================================================================ CORS
\section{Pourquoi pas de CORS ici}

Le navigateur ne parle \textbf{qu'à une seule origine}~:
\code{http://pod-royale.local}, c'est-à-dire l'arène derrière l'Ingress. La page
est rendue côté serveur par Jinja2, et \code{arena.js} appelle
\code{fetch('/api/state')} avec une URL relative, donc sur la même origine. Il
n'y a jamais de requête cross-origin, et aucun en-tête CORS n'est nécessaire.

Les guerriers et Redis ne sont \textbf{pas exposés} hors du cluster (Services
ClusterIP, pas d'Ingress). L'arène lit l'état dans Redis côté serveur, puis le
renvoie au navigateur. C'est plus simple et plus sûr~: la surface exposée se
limite à une seule application.

% ================================================================ Structure
\clearpage
\section{Structure du dépôt}

\begin{longtable}{L{2.3cm} L{13.2cm}}
\rowcolor{nuit}\entete{Dossier} & \entete{Contenu}\\
\endhead
\rowcolor{fond}\code{guerrier/} & \code{app.py}, \code{names.py}, \code{Dockerfile}\\
\code{arena/} & \code{app.py}, \code{templates/}, \code{static/}, \code{Dockerfile}\\
\rowcolor{fond}\code{k8s/} & Manifests numérotés dans l'ordre d'application
  (+ \code{bonus/}~: Gateway API)\\
\code{scripts/} & \code{up.sh} / \code{up.ps1} (tout en un), \code{build.sh},
  \code{deploy.sh}, \code{demo.sh}, \code{reset.sh}\\
\bottomrule
\end{longtable}

\end{document}
